% XMSS key generation: RFC 8391, Section 4.1.7.
% Local overlay styles keep the diagram geometry fixed between steps.
\tikzset{
  xmss hidden/.style={opacity=0,text opacity=0},
  xmss show/.code args={<#1>}{\alt<#1>{}{\pgfkeysalso{xmss hidden}}},
  xmss node/.style={circle,draw=Teal,fill=Mist,minimum size=5mm,
    inner sep=0pt}
}

% Example tree with height h=3.
\newcommand{\xmsskeygentree}{%
  \begin{tikzpicture}[x=0.68cm,y=0.68cm,font=\scriptsize,>=Latex]
    \foreach \i in {0,...,7} {
      \node[xmss node]
        (l\i) at (\i,0) {};
      % Two vertically separated keys represent the secret/public keypair.
      \begin{scope}[shift={(\i,-1.35)},draw=Gold,line width=0.7pt]
        \foreach \keyheight in {-0.16,0.16} {
          \draw (-0.15,\keyheight) circle[radius=0.075cm];
          \draw (-0.04,\keyheight) -- (0.22,\keyheight);
          \draw (0.13,\keyheight) -- (0.13,\keyheight-0.10);
          \draw (0.22,\keyheight) -- (0.22,\keyheight-0.10);
        }
      \end{scope}
      \draw[->,Teal] (\i,-1.05) --
        node[midway,right,rotate=90,anchor=south,font=\tiny] {LTree}
        (l\i.south);
    }
    \node[text=Gold] at (3.5,-1.95) {\wots keypairs};
    \foreach \i/\x in {0/0.5,1/2.5,2/4.5,3/6.5} {
      \node[xmss node] (p\i) at (\x,1) {};
      \pgfmathtruncatemacro{\leftleaf}{2*\i}
      \pgfmathtruncatemacro{\rightleaf}{2*\i+1}
      \draw[Muted] (p\i) -- (l\leftleaf);
      \draw[Muted] (p\i) -- (l\rightleaf);
    }
    \node[xmss node] (q0) at (1.5,2) {};
    \node[xmss node] (q1) at (5.5,2) {};
    \foreach \parent/\child in {q0/p0,q0/p1,q1/p2,q1/p3}
      \draw[Muted] (\parent) -- (\child);
    \node[xmss node,label=above:{public root}]
      (r) at (3.5,3) {};
    \foreach \child in {q0,q1}
      \draw[Muted] (r) -- (\child);
  \end{tikzpicture}%
}

\begin{frame}{XMSS: Overview}
  \begin{itemize}
    \setlength{\itemsep}{4mm}
    \item XMSS~\cite{DBLP:conf/pqcrypto/BuchmannDH11} is a stateful hash-based signature scheme.
    \item Security relies on the security of the underlying hash function.
    \item XMSS combines a one-time signature scheme (\wots) with a Merkle tree.
    \vspace{2mm}
    \begin{itemize}
      \setlength{\itemsep}{1mm}
      \item The root of the Merkle tree is the public key of the XMSS instance.
      \item Each leaf corresponds to a (compressed) \wots public key.
      \item A Merkle tree of height $h$ has $2^h$ leaves, and thus allows $2^h$ signing $2^h$ messages.
    \end{itemize}
  \end{itemize}
\end{frame}

% \wots key generation and signing: RFC 8391, Sections 3.1.3--3.1.5.
\begin{frame}{\wots: Key Generation}
  \small
  \begin{enumerate}
    \setlength{\itemsep}{2mm}
    \item Generate $\ell$ secret $n$-byte values
      $\mathsf{sk}_0,\ldots,\mathsf{sk}_{\ell-1}$ (randomly or pseudorandomly).
    \item Each value starts a hash chain: apply $w-1$ steps to obtain
      $\mathsf{pk}_j=\mathsf{chain}(\mathsf{sk}_j,0,w-1)$.
    \item The \wots public key is the array of all $\ell$ chain endpoints.
  \end{enumerate}
  \vspace{2mm}
  \begin{center}
    \begin{tikzpicture}[x=1.7cm,y=0.75cm,font=\scriptsize,>=Latex]
      \foreach \j/\y in {0/1,1/0,\ell-1/-1} {
        \node[anchor=east,text=Gold] at (-0.3,\y) {$\mathsf{sk}_{\j}$};
        \foreach \k in {0,...,3}
          \node[xmss node] (k\y\k) at (\k,\y) {};
        \foreach \k/\next in {0/1,1/2,2/3}
          \draw[->,Teal] (k\y\k) -- (k\y\next);
        \node[anchor=west,text=Teal] at (3.3,\y) {$\mathsf{pk}_{\j}$};
      }
      \node at (1.5,-1.65) {Example: $w=4$, three steps per chain (selected chains shown).};
    \end{tikzpicture}
  \end{center}
  {\footnotesize
    $n$: element size in bytes; $w$: Winternitz parameter;
    $\ell$: number of message and checksum digits.\\[1mm]
    Each step uses a hash key and bitmask derived from public
    $\mathsf{SEED}$ and the keypair, chain, and step address.
  }
\end{frame}

\begin{frame}{\wots: Signature Generation}
  \small
  \begin{enumerate}
    \setlength{\itemsep}{2mm}
    \item Encode the $n$-byte message digest $D$ as base-$w$ digits
      $m_0,\ldots,m_{\ell_1-1}$.
    \item Append the base-$w$ digits of the checksum
      $C=\sum_{j=0}^{\ell_1-1}(w-1-m_j)$,
      giving $\ell$ digits $a_j\in\{0,\ldots,w-1\}$.
    \item Reveal one node per chain:
      $\sigma_j=\mathsf{chain}(\mathsf{sk}_j,0,a_j)$.
      The signature is $(\sigma_0,\ldots,\sigma_{\ell-1})$.
  \end{enumerate}
  \vspace{2mm}
  \begin{center}
    \begin{tikzpicture}[x=1.7cm,y=0.75cm,font=\scriptsize,>=Latex]
      \foreach \j/\y/\digit in {0/1/1,1/0/3,2/-1/2} {
        \node[anchor=east] at (-0.3,\y) {$a_{\j}=\digit$};
        \foreach \k in {0,...,3}
          \node[xmss node,draw=Muted,fill=white] (s\j\k) at (\k,\y) {};
        \foreach \k/\next in {0/1,1/2,2/3} {
          \ifnum\k<\digit
            \draw[->,Teal,thick] (s\j\k) -- (s\j\next);
          \else
            \draw[->,Muted,dashed] (s\j\k) -- (s\j\next);
          \fi
        }
        \node[xmss node,draw=Gold,fill=Gold!20] at (s\j\digit) {};
        \node[anchor=west,text=Gold] at (3.3,\y) {$\sigma_{\j}$: step $\digit$};
      }
      \node at (1.5,-1.65) {Example: $w=4$; gold nodes form the signature (selected chains shown).};
    \end{tikzpicture}
  \end{center}
  {\footnotesize
    Verification completes each chain with $w-1-a_j$ steps to recover its endpoint.\\[1mm]
    \alert{Each \wots keypair signs at most one message digest.}
  }
\end{frame}

\begin{frame}{XMSS: Key Generation}
  \small
  \setbeamerfont{itemize/enumerate body}{size=\small}
  \begin{columns}[T,onlytextwidth]
    \begin{column}{0.54\textwidth}
      \begin{enumerate}
        \item Start from $2^h$ \wots keypairs: one for each leaf
          of the tree.
        \item Compress each \wots public key via an L-tree
          (LTree) to an $n$-byte value. These values are the leaves.
        \item Hash pairs of nodes upward, level by level,
          to obtain the $n$-byte root.
      \end{enumerate}
      {\setbeamerfont{block title}{size=\small}
        \setbeamerfont{block body}{size=\footnotesize}
        \begin{block}{Scheme parameters}
          $h$: tree height ($2^h$ leaves).\\
          $n$: hash output length in bytes.
        \end{block}
}
    \end{column}
    \begin{column}{0.43\textwidth}
      \centering
      \vspace{3mm}
      \xmsskeygentree\\[3mm]
      \small Example: $h=3$

    \end{column}
  \end{columns}
  \vspace{0mm}
  {\footnotesize
    \textbf{Public key:} $\mathsf{PK}=(\mathsf{root},\mathsf{SEED})$, with public seed $\mathsf{SEED}$.\\
    \textbf{Secret key:} $\mathsf{SK}=(i,\{\mathsf{sk}^{\mathrm{\wots}}_j\}_{j=0}^{2^h-1},\mathsf{SK}_{\mathrm{PRF}},\mathsf{PK})$.\\
    $\mathsf{SK}_{\mathrm{PRF}}$: secret key for message randomization.
  }
\end{frame}

\begin{frame}{\xmss: Signature Verification}
  \relax
  {\small
    \textbf{Inputs}: message $M$, $\mathsf{PK}=(\mathsf{root},\mathsf{SEED})$,
    $\sigma=(i,R,\sigma_{\mathrm{\wots}},\mathsf{auth}_i)$, with $0\leq i<2^h$.\par}
  \vspace{2mm}
  \begin{columns}[T,onlytextwidth]
    \begin{column}{0.54\textwidth}
      \begin{enumerate}
        \small
        \setlength{\itemsep}{0.5mm}
        \item<1-> Recompute the message digest $D = H_{\mathrm{msg}}(R\mathbin\Vert\mathsf{root}\mathbin\Vert i,M)$.
        \item<2-> Complete the \wots chains from
          $\sigma_{\mathrm{\wots}}$ to recover $\mathsf{pk}_i$.
        \item<3->LTree compression gives leaf $i$.
        \item<4-> Hash upward with the $h$ authentication siblings.
          Bits of $i$ determine left/right order.
        \item<5-> Accept if the reconstructed root equals
          $\mathsf{root}$ in the public key.
      \end{enumerate}
    \end{column}
    \begin{column}{0.43\textwidth}
      \centering
      \vspace{2mm}
      \begin{tikzpicture}[x=0.68cm,y=0.7cm,font=\scriptsize]
        \foreach \i in {0,...,7}
          \node[xmss node,draw=Muted!45,fill=white] (v\i) at (\i,0) {};
        \foreach \i/\x in {0/0.5,1/2.5,2/4.5,3/6.5} {
          \node[xmss node,draw=Muted!45,fill=white] (p\i) at (\x,1) {};
          \pgfmathtruncatemacro{\leftleaf}{2*\i}
          \pgfmathtruncatemacro{\rightleaf}{2*\i+1}
          \draw[Muted!40] (p\i) -- (v\leftleaf);
          \draw[Muted!40] (p\i) -- (v\rightleaf);
        }
        \node[xmss node,draw=Muted!45,fill=white] (q0) at (1.5,2) {};
        \node[xmss node,draw=Muted!45,fill=white] (q1) at (5.5,2) {};
        \node[xmss node,draw=Muted!45,fill=white,label=above:{public root}]
          (r) at (3.5,3) {};
        \foreach \parent/\child in {q0/p0,q0/p1,q1/p2,q1/p3,r/q0,r/q1}
          \draw[Muted!40] (\parent) -- (\child);
        \foreach \name in {v3,p0,q1}
          \node[xmss node,draw=Gold,fill=Gold!20] at (\name) {};
        \foreach \name/\step in {v2/2,p1/3,q0/4,r/5}
          \node[xmss node,fill=Teal!25,xmss show=<\step->] at (\name) {};
        \foreach \parent/\child/\step in {p1/v2/3,q0/p1/4,r/q0/5}
          \draw[Teal,line width=1.2pt,xmss show=<\step->]
            (\parent) -- (\child);
        \node[below=2mm of v2,align=center,xmss show=<2->]
          {$\mathrm{LTree}(\mathsf{pk}_i)$\\$i=2$};
      \end{tikzpicture}\\[1mm]
      {\footnotesize\color{Gold}Gold: supplied authentication siblings.}\\
      {\footnotesize\color{Teal}Teal: reconstructed path.}
    \end{column}
  \end{columns}
\end{frame}

\begin{frame}{\xmssmt: Multi-Tree Variant}
  \begin{columns}[T,onlytextwidth]
    \begin{column}{0.53\textwidth}
      \begin{itemize}
        \small
        \setlength{\itemsep}{2mm}
        \item<1-> Split total height $h$ into $d$ layers of \xmss
          trees, each of height $h/d$.
        \item<2-> Bottom-layer \wots keys sign message digests.
        \item<3-> Upper-layer \wots keys sign the roots of trees
          in the layer below.
        \item<4-> Verify layer by layer up to the single public root.
      \end{itemize}
      \uncover<4->{%
        {\setbeamerfont{block title}{size=\small}
          \setbeamerfont{block body}{size=\footnotesize}
          \begin{block}{Tradeoff}
            $2^h$ messages; smaller trees to generate.\\[1mm]
            Signature: $d$ \wots signatures and $h$ path nodes.\\[1mm]
            For $d=1$, \xmssmt is functionally equivalent to \xmss.
          \end{block}}}
    \end{column}
    \begin{column}{0.44\textwidth}
      \centering
      \vspace{2mm}
      \begin{tikzpicture}[x=1cm,y=0.75cm,font=\scriptsize,>=Latex]
        \node[xmss node,label=above:{public root}] (r) at (2.4,4.6) {};
        \node[xmss node] (p0) at (1.2,3.9) {};
        \node[xmss node] (p1) at (3.6,3.9) {};
        \foreach \parent in {p0,p1}
          \draw[Teal] (r) -- (\parent);
        \foreach \j/\x/\parent in {0/0.6/p0,1/1.8/p0,2/3.0/p1,3/4.2/p1} {
          \node[xmss node] (u\j) at (\x,3.1) {};
          \draw[Teal] (\parent) -- (u\j);
          \coordinate (b\j) at (\x,1.9);
          \draw[Teal,fill=Mist,xmss show=<2->]
            (b\j) -- (\x-0.5,0.9) -- (\x+0.5,0.9) -- cycle;
          \draw[->,Gold,thick,xmss show=<3->] (b\j) -- (u\j.south);
        }
        \node[xmss node,xmss show=<2->] (leaf) at (1.8,0.9) {};
        \node[xmss show=<2->] (digest) at (1.8,0.1) {$D$};
        \draw[->,Gold,thick,xmss show=<2->] (digest) -- (leaf);
        \node[text=Muted,anchor=east] at (4.7,4.7) {top layer};
        \node[text=Muted,xmss show=<2->,anchor=east] at (4.7,0.2)
          {bottom layer};
        \node[text=Gold,xmss show=<3->,fill=white,inner sep=1pt]
          at (2.4,2.5) {\wots};
      \end{tikzpicture}\\[2mm]
      {\footnotesize Example: $h=4$, $d=2$.}\\[1mm]
      {\scriptsize\color{Gold}Gold arrows: \wots signs a digest or root.}
    \end{column}
  \end{columns}
  % RFC 8391, Section 4.2. Upper-layer signatures can be reused for the
  % same lower-tree root; a \wots key must never sign distinct roots.
\end{frame}
